\subsubsection*{Реализация серверной части}

Серверная часть разработана на базе ASP.NET Core 9 как монолитное веб-приложение, внутри которого реализованы логически отделённые компоненты для управления различными аспектами системы онбординга. Архитектура предусматривает модульность кода через разделение ответственности между сервисами, однако они развёртываются в едином приложении, а не как отдельные микросервисы. Во время преддипломной практики основное внимание было уделено разработке следующих компонентов:

\textbf{1. Компонент отслеживания прогресса и геймификации:}

Сервис, реализованный в классе `GamificationService`, отвечает за всю логику геймификации, аналитики и отслеживания прогресса пользователей:
\begin{itemize}
	\item \textbf{Система XP и уровней:} 
	\begin{itemize}
		\item Алгоритм расчёта опыта: +50 XP за завершение модуля, +100 XP за успешное прохождение теста, +10 XP за ежедневную активность.
		\item Автоматическое определение уровня на основе накопленного XP (100 XP за уровень, максимум 50 уровней).
		\item Сохранение истории всех начисления XP в БД для аналитики и аудита.
	\end{itemize}
	
	\item \textbf{Достижения:} Реализована система из 12 типов достижений с автоматической проверкой условий получения и уведомлением пользователя.
	
	\item \textbf{Leaderboard (Таблица лидеров):}
	\begin{itemize}
		\item Endpoint GET /api/leaderboard для получения рейтинга с сортировкой по XP.
		\item Кэширование в Redis с TTL 5 минут для обеспечения быстрого доступа.
		\item Обновления рейтинга в реальном времени через SignalR WebSocket при изменении позиций.
		\item Поддержка фильтрации по отделам и корпоративным иерархиям.
	\end{itemize}
	
	\item \textbf{Аналитика:} 
	\begin{itemize}
		\item Расчёт ключевых метрик: средний XP по организации, процент завершения модулей, активность пользователей по времени.
		\item Таблица ActionLogs с записью всех действий пользователей (просмотры модулей, прохождение тестов, получение достижений) для детального анализа.
		\item Прогнозирование времени завершения онбординга на основе текущего прогресса и среднего времени прохождения.
		\item Endpoints: /api/analytics/metrics, /api/analytics/export.
	\end{itemize}
	
	\item \textbf{Экспорт данных:} Генерация отчётов в PDF (QuestPDF) и Excel (ClosedXML) с полной статистикой по пользователям и отделам.
\end{itemize}

\textbf{2. Компонент интеллектуального помощника (AI Mentor):}

Сервис в классе `AiMentorService` обеспечивает интеграцию с внешним API обработки естественного языка для предоставления ответов на вопросы пользователей:
\begin{itemize}
	\item \textbf{Chat Endpoint:} POST /api/ai/chat для отправки вопросов и получения ответов от AI.
	\begin{itemize}
		\item Поддержка потокового ответа через Server-Sent Events (SSE) для улучшенного UX.
		\item Использование модели LLaMA-2 от Groq для обработки естественного языка.
	\end{itemize}
	
	\item \textbf{Кэширование:}
	\begin{itemize}
		\item Redis кэш для часто задаваемых вопросов с TTL 30 дней.
		\item Сокращение времени отклика на 60\% благодаря кэшированию.
		\item Автоматическая инвалидация кэша при получении новых вопросов.
	\end{itemize}
	
	\item \textbf{Хранение истории:} Сохранение всех диалогов в БД с метаинформацией (пользователь, время, время отклика).
	
	\item \textbf{Мониторинг и логирование:} Логирование всех запросов и ответов для отладки, улучшения качества и анализа популярных вопросов.
	
	\item \textbf{AI-рекомендации:} На основе анализа истории взаимодействия пользователя система предлагает:
	\begin{itemize}
		\item Рекомендуемые модули обучения.
		\item Подходящих наставников на основе их опыта и специализации.
		\item Уведомления о необходимых действиях для ускорения прогресса.
	\end{itemize}
\end{itemize}

\textbf{3. Компонент уведомлений (Notification Service):}

Реализован в классе `NotificationService` с поддержкой real-time и email-уведомлений:
\begin{itemize}
	\item SignalR Hub для real-time обновлений: таблица лидеров, достижения, сообщения наставника и статус онбординга.
	\item Email-уведомления через SMTP с шаблонами для завершения модулей, достижений, назначений наставников и еженедельных дайджестов.
	\item Хранение истории уведомлений с возможностью фильтрации, поиска и удаления.
	\item Управление предпочтениями уведомлений на уровне пользователя.
\end{itemize}

\textbf{4. Компонент интеграции с системой RIMS (RimsIntegrationService):}

Сервис обеспечивает синхронизацию кадровых данных из корпоративной системы управления персоналом:
\begin{itemize}
	\item Синхронизация пользователей из RIMS: ФИО, должность, отдел, email.
	\item Поддержка организационной иерархии для распределения обучения по отделам.
	\item Обновление статусов сотрудников (новичок, активный, уволен) для управления доступом и обучением.
	\item Endpoint GET /api/rims/sync и автоматический scheduler на 02:00 ежедневно.
	\item Логирование операций синхронизации для аудита и отладки.
\end{itemize}

\textbf{5. Дополнительные компоненты:}

Приложение включает вспомогательные компоненты для управления аутентификацией, авторизацией и контентом, интегрированные в единое монолитное приложение:

\textbf{Управление аутентификацией и авторизацией:} Реализовано в сервисах `AuthenticationProvider`, `JwtTokenGenerator` и `AuthorizationService` для управления пользователями, ролями и JWT-токенами (access-токены на 24ч, refresh-токены на 7 дней) с интеграцией синхронизации персонала из RIMS.

\textbf{Управление контентом:} Контроллеры и сервисы для управления модулями обучения, контентом, системой тестирования и материалами.

\textbf{Сервисы экспорта и отчётности:} Класс `ReportExportService` для генерации отчётов в различных форматах.

\textbf{Инфраструктура приложения:}
\begin{itemize}
	\item \textbf{SQL Server:} Реляционная БД для основного хранилища данных с поддержкой таблиц для пользователей, модулей, прогресса, логов действий и других сущностей.
	\item \textbf{Redis (опционально):} Возможность использования кэширования для оптимизации производительности (кэширование leaderboard, AI ответов, сессий).
	\item \textbf{SignalR:} Интеграция для real-time WebSocket соединений в рамках одного приложения для обновления таблицы лидеров, уведомлений о достижениях и статуса онбординга.
	\item \textbf{Entity Framework Core:} ORM для взаимодействия с БД, управления миграциями и LINQ-запросами.
	\item \textbf{Архитектура приложения:} Все компоненты зарегистрированы и управляются через ASP.NET Core Dependency Injection контейнер в едином процессе приложения.
\end{itemize}

\textbf{Внешние интеграции:}
\begin{itemize}
	\item \textbf{Внешние API для обработки естественного языка:} Доступ к LLM-моделям через API для обработки вопросов в AI Chat Assistant.
	\item \textbf{RIMS API:} Синхронизация пользователей и структуры организации из корпоративной системы управления персоналом.
	\item \textbf{Email (SMTP):} Отправка уведомлений пользователям через шаблоны писем.
\end{itemize}

\textbf{Стек технологий серверной части:} ASP.NET Core 9, Entity Framework Core, SQL Server, SignalR, xUnit (тестирование), Docker, интеграция с внешними API (Groq, RIMS, SMTP для уведомлений).

\textbf{Общий объём разработанного кода:} Более 25 000 строк на backend, включая бизнес-логику, валидацию, обработку ошибок, unit-тесты и integration-тесты.
